% Source: https://artofproblemsolving.com/community/c3h2419202_silly_doubt_from_number_theory_marathon
% Stable identity: aops-topic-2419202-q01
\textbf{Why modulo $100$.} Any positive integer $N$ can be written as $N=100q+r$ with $0\le r\le99$, where $r$ is the number formed by the last two digits of $N$. So the tens digit of $N$ is the tens digit of the remainder $r=N\bmod100$, and we only need $N$ modulo $100$.

\textbf{The computation.} Let $e=10111213141516171819$. Since $123456789\equiv89\equiv-11\pmod{100}$ and $e$ is odd,
\[123456789^{e}\equiv(-11)^e=-11^e\pmod{100}.\]
By the binomial theorem,
\[11^e=(1+10)^e=1+\dbinom e1\cdot10+\dbinom e2\cdot10^2+\cdots\equiv1+10e\pmod{100}.\]
As $e\equiv9\pmod{10}$, we have $10e\equiv90\pmod{100}$, so $11^e\equiv91$ and
\[123456789^{e}\equiv-91\equiv9\pmod{100}.\]
The number ends in the digits $09$, so its tens digit is $\boxed{0}$.
